% ECONOMICS_PROBLEM: {"id": "C78-225", "title": "Maximum number of reachable deferred-acceptance states", "jel_code": "C78", "source_id": "OP-0205"}
% !TeX program = xelatex
\documentclass[11pt,a4paper]{article}
\usepackage[margin=23mm]{geometry}
\usepackage{fontspec}
\setmainfont{Latin Modern Roman}
\usepackage{amsmath,amssymb,mathtools}
\usepackage{xcolor,enumitem,booktabs,longtable,array,xurl,hyperref}
\definecolor{muted}{HTML}{53636C}
\hypersetup{colorlinks=true,urlcolor=blue,linkcolor=blue}
\setlength{\parindent}{0pt}
\setlength{\parskip}{5pt}
\newcommand{\R}{\mathbb R}
\newcommand{\E}{\mathbb E}
\newcommand{\N}{\mathbb N}
\newcommand{\ind}{\mathbf 1}
\newcommand{\Var}{\operatorname{Var}}
\newcommand{\Cov}{\operatorname{Cov}}
\newcommand{\supp}{\operatorname{supp}}
\DeclareMathOperator*{\argmin}{arg\,min}
\DeclareMathOperator*{\argmax}{arg\,max}
\newcommand{\entrymeta}[3]{{\sffamily\small\color{muted}\textbf{\texttt{#1}}\quad|\quad #2\par #3\par}}
\newcommand{\Problemref}[1]{\texttt{#1}}
\newcommand{\JELref}[2]{\texttt{#2} (JEL \texttt{#1})}
\begin{document}
\section*{Maximum number of reachable deferred-acceptance states}
\textbf{Problem ID:} \texttt{C78-225}\quad \textbf{JEL:} \texttt{C78}\par
% BEGIN ECONOMICS STATEMENT
\label{problem:OP-0205}
% GAME_THEORY_PROBLEM: {"local_id": "GT-075", "original_number": 75, "kind": "P", "entry_kind": "statement_only", "published": false, "review_required": true, "category": "game_theory"}
\label{game:GT-075}
{\sffamily\small\color{muted}
\noindent\textbf{ID: GT-075.}\quad\textbf{Category:} Extremal matching combinatorics.
\quad\textbf{Kind: P.}\quad\textbf{Status:} Open extremal sequence in the cited source.
\par}
\subsection*{Definitions and mathematical statement}
Let \(\mathcal I_n\) be all strict complete preference profiles on fixed
disjoint labeled sets \(P,R\) with \(|P|=|R|=n\).
For \(I\in\mathcal I_n\), start deferred acceptance with no held proposer
and counters \(q(p)=0\). A state is a pair \((h,q)\), where
\(h:R\to P\cup\{\bot\}\) records distinct held proposers and
\(q(p)\in\{0,\ldots,n\}\) records the number of proposals made by \(p\).
At a legal step an unheld proposer with \(q(p)<n\) proposes to its
next untried receiver. Its counter increases by one and the receiver
keeps its preferred of the new proposer and its previous holder;
a rejected proposer is unheld.
Define \(\operatorname{Reach}(I)\) as all states reached from the
initial state by all finite legal serial schedules, including the
initial and terminal states, and define
\[
 N_{\max}(n)=
       \max_{I\in\mathcal I_n}|\operatorname{Reach}(I)|\qquad(n\geq 1).
\]
The maximum exists because \(\mathcal I_n\) is finite.

\paragraph{Problem.}
Determine the extremal sequence \(N_{\max}(n)\), by an exact formula
or a sharp structural characterization valid for all \(n\).

\subsection*{Short English statement}
How many different execution states can an order-\(n\) stable-marriage
instance have?

\subsection*{Variants and scope boundaries}
No primality restriction is imposed. This counts states, not the number
of schedules, proposals, or stable matchings. Exhaustive enumeration is
a finite algorithm for each fixed \(n\), but does not resolve the requested
general extremal characterization.

\subsection*{Source relationship and status}
This is [S1, Section 7, problem B].
The source reports exact values \(2,5,15,52\) for \(n=1,2,3,4\),
and only \(N_{\max}(5)\geq 232\). Thus a Bell-number continuation of
the first four values is already impossible.

\subsection*{References}
\begingroup\footnotesize\raggedright
\noindent[S1] Yoshiteru Ishida,
\emph{Prime Factorisation of Stable-Matching Instances: Uniqueness,
Simultaneous Products, and an Exact Census}, 2026,
arXiv:2610.05617v1.
\url{https://arxiv.org/html/2610.05617v1}.
Locators: Section 2 (reachable-state definition);
Proposition 21.4, final discussion; Section 7, problem B.
\par\endgroup

\subsection*{Common conventions: Source mathematical conventions}

\textbf{Finite games.} For a finite set $A$, $\Delta(A)$ is its probability
simplex. Players choose independent mixed strategies in Nash equilibrium;
correlation is allowed only when explicitly part of the solution concept.
In a game with payoff functions $u_i$ and strategy profile $x$, an additive
$\varepsilon$ Nash equilibrium satisfies
\[
 u_i(x)\geq u_i(x_i',x_{-i})-\varepsilon
 \quad\text{for every player $i$ and every admissible deviation $x_i'$}.
\]
For numerical approximation, payoffs are normalized as locally specified.
An $\varepsilon$ payoff residual is distinct from distance at most
$\varepsilon$ to the set of exact equilibria. Point convergence, convergence
of empirical play, and convergence in distance to an equilibrium set are
also distinct.

\textbf{Computational representations.} Unless stated otherwise, a complexity
question takes rational data with binary encoding; polynomial time is measured
in total bit length. A fixed-accuracy polynomial algorithm is not necessarily
polynomial in $1/\varepsilon$, and neither is necessarily polynomial in
$\log(1/\varepsilon)$. Succinct games and value, demand, or sampling oracles
must declare their interfaces. Exact real solutions require an explicit
representation; an unspecified list of arbitrary real numbers is not a
finite computational output. A claimed algorithmic barrier is meaningful
only for its specified model and reductions.

\textbf{Dynamic games.} Histories, information, observation, and allowable
randomization are part of the model. A stationary strategy depends only on
the current state; a positional strategy in a deterministic graph game
chooses one action at each controlled vertex. Nash equilibrium at an initial
state and equilibrium after every history are different requirements.
An expectation of a pathwise $\liminf$ payoff, a $\liminf$ of expected averages,
a limiting expected average, and a uniform finite-horizon equilibrium are
different criteria. None is silently substituted for another.

\textbf{Allocation and mechanisms.} A complete allocation assigns every item
exactly once unless otherwise stated. Zero-valued goods, negative values,
capacity constraints, feasibility, lotteries, and copies of goods affect
fairness definitions and are specified locally. Dominant-strategy incentive
compatibility, Bayesian incentive compatibility, universal truthfulness,
and truthfulness in expectation impose different inequalities. Whenever
an objective ratio has zero denominator, its convention is stated or those
instances are excluded.

\textbf{Infinite-dimensional models.} Expectations, integrals, stochastic
equations, and optimization objectives require their local regularity,
integrability, filtrations, and admissible domains. A backward--forward
mean-field system is a boundary-value problem; it is not an initial-value
semiflow without an additional construction. Long-time, large-population,
and vanishing-noise limits require an order of limits and a convergence
topology. If a minimum need not be attained, the objective is its infimum
or supremum and the approximation requirement is stated separately.
% END ECONOMICS STATEMENT
\end{document}
